% =======================================================================
\section{Introduction}

\subsection{Background}

Competitive collective systems are agents distributed in a bounded
domain, coordinating internally while responding to an opponent.
Autonomous vehicle fleets sharing road space, opposing crowd flows at
transport hubs, predator--prey populations in an ecosystem, and sports
teams contesting territory all share this structure. Consecutive
positions are coupled by that opposition (competitive dependence); this
paper asks only whether homology can be attributed by organisational
scale. Such systems operate across multiple spatial scales at once:
individual agent decisions at the scale of metres, small-group
coordination at tens of metres, and whole-system organisation at
hundreds of metres. Understanding how these scales interact is a
standing challenge in applied mathematics, and one that high-frequency
spatial tracking data is increasingly able to inform
\citep{gudmundsson2017}.

Agents within such systems can arrange themselves in ring-like
formations that enclose empty space: a pressing trap encircling a player in possession, or a gap between two organisational units that creates room to manoeuvre. These arrangements have a topological
character: they are loops in the position data, and their persistence
over time distinguishes deliberate structural organisation from
transient clustering. Conventional geometric summaries capture different aspects of spatial
dispersion \citep{gudmundsson2017}. Convex-hull area, nearest-neighbour
distance, and the range of agent positions along each axis do not
directly quantify loop structure.

Persistent homology provides a principled way to detect and measure such
formations across scale \citep{edelsbrunner2010,carlsson2009}. A
filtration of simplicial complexes indexed by a scale parameter tracks
the birth and death of topological features. Each feature is recorded as
a bar: the interval from the scale at which it appears (birth) to the
scale at which it disappears (death). The collection of bars is the
persistence diagram \citep{cohensteiner2007}.
Figure~\ref{fig:pipeline}(c) shows a schematic of finite $H_1$ pairs.

That representation has recovered cavities in protein atom clouds that
mark flexibility and folding \citep{xia2014}. It has identified
medium-range order in amorphous solids that conventional crystallography
does not resolve \citep{hiraoka2016}. In biological aggregation models
it has shown that distinct swarming organisations appear at different
filtration values \citep{topaz2015}. The same single-parameter
construction, applied directly to a multi-agent point cloud, raises the
measurement difficulty treated in Section~\ref{sec:scaleconflation}.

\subsection{The Scale Conflation Problem}
\label{sec:scaleconflation}

A single-parameter Vietoris--Rips filtration on the full point cloud
records features at many filtration values in one persistence diagram.
Bars born at different distances therefore occupy the same object,
whether they arise from local proximity, small-group arrangement, or
system-wide organisation.

That mixing is scale conflation, and it is a property of the filtration
on a hierarchically organised point cloud, not of competition. Birth
records when a loop closes among the original points. It does not name
the organisational level of the units that enclose the loop. In a
competitive collective system the hierarchy has named levels
(individual, small group, envelope), so a superimposed diagram cannot
attribute an $H_1$ feature to a named level, and it cannot track that
level through time. The metre values of those levels are not given in
advance; they are recovered by a cutoff sweep
(Section~\ref{sec:cutoffselection}).
Aggregation models already show distinct organisational structures at
different filtration values \citep{topaz2015}. Multi-agent snapshots
likewise yield Vietoris--Rips summaries that depend on the chosen scale
\citep{gu2022}.

Two existing responses address related difficulties. Multiparameter
persistence treats several scales as simultaneous algebraic parameters
\citep{botnan2022}. Computational cost currently limits that machinery
to point clouds much smaller than the 22-agent, high-frequency case
used here. Persistent homology of multiscale clusterings analyses the
hierarchy produced by the clustering itself \citep{schindler2023}. It
does not yield a scale-attributable diagram on a reduced point cloud.
Section~\ref{sec:relatedwork} reviews these lines in more detail.

Organisational level is assigned here by changing the input, not by
slicing bars of one diagram. Hierarchical clustering at validated cutoff
distances decomposes the cloud; persistent homology is then computed on
the centroids. Clustering changes the geometry. A Vietoris--Rips range
suited to one cutoff can return no $H_1$ at another. A data-driven
truncation of the filtration parameter removes that obstruction
(Section~\ref{sec:adaptivefiltration}).

\subsection{Contributions}

The measurement pipeline has two steps. Domain-informed hierarchical
clustering, with cutoffs checked against independent clustering-quality
metrics, separates organisational levels before homology is computed. A
data-driven truncation of the Vietoris--Rips parameter keeps $H_1$
detection from collapsing once clustering has changed the geometry.
Neither step is algebraically new
\citep{carlsson2010clustering,schindler2023}; the content is that the
combination survives broadcast tracking.

The empirical findings are four. A sweep of the cutoff over
$[0.5, 30.0]$~m identifies three stable $H_0$ regimes across ten
matches, not a single demonstration case. $H_1$ is detected at two of
the three scales, with a geometric cycle representative for every
feature in the primary-match analysis. The two $H_1$ scales carry
complementary rather than redundant information. Event association is
reported only as construct validity in Section~\ref{sec:eventcorr}.

\subsection{Related Work}
\label{sec:relatedwork}

Persistent homology was introduced by \citet{edelsbrunner2002} and
placed on firm computational foundations by \citet{zomorodian2005}.
Stability of persistence diagrams under perturbation was established by
\citet{cohensteiner2007}, and efficient computation via the
Vietoris--Rips filtration is provided by \citet{bauer2021ripser}.
Persistence landscapes provide a functional representation suitable for
statistical analysis \citep{bubenik2015}.

Applications of topological data analysis (TDA) to collective behaviour
include
\citet{topaz2015}, who applied the method to biological aggregation
models and visualised scale-dependent structure by plotting Betti
numbers across both simulation time and filtration value.
\citet{gu2022} applied Vietoris--Rips persistent homology to multi-agent
system snapshots for change-point detection, demonstrating the utility
of topological features for temporal dynamics, though their analysis
operates at a single scale rather than decomposing by organisational
level. The multiparameter persistence literature, surveyed by
\citet{botnan2022}, provides a theoretical framework for genuinely
simultaneous multi-scale analysis, but computational cost currently
limits its application to point clouds much smaller than the 22-agent,
high-frequency case considered here. \citet{schindler2023} address
hierarchical decomposition directly, developing tools for analysing
multiscale clusterings with persistent homology. The present work asks
whether cluster-then-persistent homology, with the data-driven truncation
of Section~\ref{sec:adaptivefiltration}, survives a real, high-frequency,
noisy point cloud and yields stable scale regimes across ten independent
matches.

In sports analytics specifically, spatial analysis methods including
pitch-control models and spatio-temporal pattern recognition are
surveyed in \citet{gudmundsson2017}. Related perspectives model teams as
coordinated multi-agent systems and apply network metrics to passing and
collective structure \citep{buldu2018,grund2012}; these approaches
complement the present pipeline but do not quantify two-dimensional
enclosure through persistent homology. To our knowledge, no prior work
combines domain-informed scale decomposition with persistent homology
for a competitive collective system at high frequency, or validates the
resulting scale regimes across multiple independent matches.
